\section{Justificación}
\estado{EN CONSTRUCCIÓN}

\subsection{Justificación}
Pregunta de investigación: 

¿Cómo afectan los errores de estimación del tiempo de viaje de Google Maps a los estudiantes de la UPTC en Tunja que viajan al menos 3 veces por semana?

Este estudio busca medir y describir cómo afectan a los estudiantes de la UPTC los errores de estimación del tiempo de viaje de Google Maps. Para eso compara el tiempo que la aplicación estima con el tiempo real de los trayectos, y analiza qué consecuencias tiene esa diferencia en su rutina académica, personal y económica. En el problema ya identificado, este estudio cumple el papel de pasar de la sospecha a la evidencia. Hoy se sabe que la aplicación se usa mucho, pero no hay datos propios sobre qué tan bien funciona en Tunja ni sobre cómo influyen sus errores en quienes dependen de ella todos los días. Lo que se espera obtener es justamente ese conocimiento: qué tan grande es el error y qué efectos tiene en los estudiantes.

Es necesario hacerlo ahora porque la dependencia de las aplicaciones de navegación crece cada año \cite{dahmani2020,elpais2026} y estas ofrecen tiempos estimados que los usuarios usan para planear sus trayectos \cite{bastarianto2025}. Aun así, esas estimaciones pueden fallar, y en otros estudios se han encontrado diferencias entre el tiempo que calculan las aplicaciones y el tiempo que reportan los viajeros \cite{bastarianto2025,arxiv2017}. A esto se suma que las condiciones viales de Tunja, con congestión en horas pico \cite{uptc2022} y un transporte colectivo con deficiencias \cite{sanchez2012}, hacen dudoso que las estimaciones funcionen igual de bien que en otras ciudades. Los estudiantes de la UPTC, además, tienen horarios fijos y se desplazan al menos tres veces por semana, así que un error pequeño se repite y se acumula con el tiempo. Mientras no exista evidencia, cada semestre pasa con estudiantes llegando tarde, perdiendo clases o evaluaciones y gastando más en transporte, sin que nadie sepa si la causa es la estimación de la aplicación.

El estudio sirve, en primer lugar, a los propios estudiantes, porque les permite saber cuánto margen de tiempo conviene dejar al confiar en la aplicación y así evitar llegadas tarde, pérdida de clases y gastos innecesarios. También sirve como base para otros estudios sobre movilidad y aplicaciones de navegación en ciudades intermedias como Tunja, donde hay poca información propia. Su aporte está en cubrir tres dimensiones a la vez (académica, personal y económica) en lugar de limitarse a un solo efecto como la puntualidad.

El estudio se aborda comparando dos cosas: el tiempo estimado por la aplicación y el tiempo real medido en los trayectos, junto con la percepción de los estudiantes sobre las consecuencias de esa diferencia. Este enfoque es más conveniente que una encuesta de percepción aislada, porque permite relacionar el error con datos medibles y no solo con opiniones, y así saber si lo que los estudiantes sienten corresponde a lo que pasa realmente. Además, deja datos organizados y una forma de medir que se pueden repetir en otras rutas, otras aplicaciones u otras ciudades.

\notaseccion{Se redacta en la Fase 3 del Curso (Formulación de la propuesta).}
